\bookchapter{บทที่ 48: ทำความรู้จัก Android—จากระบบปฏิบัติการสู่ระบบนิเวศนักพัฒนา}

\chapterintro{ยินดีด้วยที่คุณตัดสินใจก้าวนี้—จาก “คนเล่นมือถือ” มาเป็น “คนทำแอปมือถือ” ก่อนลงมือ มาแผ่แผนที่ดูกันก่อน: Android มาจากไหน? ทำไมนักพัฒนาทั่วโลกถึงแห่กันเขียนแอปให้มัน? จากนิ้วที่พิมพ์คีย์บอร์ดของคุณจนถึงโทรศัพท์ผู้ใช้ แอปหนึ่งตัวเดินผ่านอะไรมาบ้าง?}

\booksection{ชาติกำเนิดของ Android: จาก “หุ่นยนต์ตัวจิ๋ว” สู่ “ระบบที่ใหญ่ที่สุดในโลก”}

เรื่องราวของ Android เริ่มปี 2003—วิศวกรชื่อ Andy Rubin ก่อตั้งบริษัท Android ขึ้น เป้าหมายเรียบง่าย: ทำระบบปฏิบัติการอัจฉริยะที่ “เข้าใจกล้องมากกว่าใคร” ต่อมา Google เข้าซื้อกิจการในปี 2005 และปี 2007 ร่วมกับผู้ผลิตหลายสิบรายก่อตั้ง Open Handset Alliance (OHA) จนเดือนกันยายน 2008 โทรศัพท์ Android รุ่นแรก HTC Dream (G1) วางขายอย่างเป็นทางการ

เรื่องต่อจากนั้นทุกคนก็รู้แล้ว: ถึงปี 2026 Android คือระบบปฏิบัติการเคลื่อนที่ที่มีส่วนแบ่งตลาดมากที่สุดในโลก อุปกรณ์ที่ใช้งานอยู่มีจำนวนหลักพันล้านเครื่อง โทรศัพท์ของคุณ กล่องทีวี จอใหญ่ในรถ นาฬิกาอัจฉริยะ ส่วนใหญ่อาจวิ่งบน Android ทั้งหมด นั่นแปลว่าอะไร? แปลว่าโค้ดทุกบรรทัดที่คุณเขียน มีโอกาสถูกใช้งานโดยคนหมู่มากที่สุดในโลก

\begin{tcbinfo}{เคล็ดลับ: พันธุกรรมโอเพนซอร์สของ Android}
โค้ดแกนของ Android (AOSP, Android Open Source Project) เป็นโอเพนซอร์ส—ผู้ผลิตรายใดก็เอาไปแก้และปรับแต่งได้ เพราะงั้น Xiaomi จึงมี MIUI (ตอนนี้ชื่อ HyperOS), Huawei มีชั้นรองรับความเข้ากันได้ของ HarmonyOS, ส่วน Samsung มี One UI อันนี้ดีต่อนักพัฒนา: ที่คุณเรียนคือ “การพัฒนา Android” ไม่ใช่ “การพัฒนาของแบรนด์ใดแบรนด์หนึ่ง” ฝีมือชุดเดียว ใช้ได้กับทุกแบรนด์

\end{tcbinfo}

การวิวัฒนาการของเลขเวอร์ชันก็สนุกไม่แพ้กัน: เริ่มจาก 1.5 ชื่อ Cupcake ทุกเวอร์ชันใหญ่ใช้ชื่อขนมหวาน—Donut, Eclair, Froyo, Gingerbread... จนถึง Android 10 ปี 2019 จึงเปลี่ยนมาใช้ตัวเลขล้วน ถึงเดือนสิงหาคม 2026 เวอร์ชันทางการล่าสุดคือ Android 17 ส่วนตอนพัฒนา “เวอร์ชันขั้นต่ำที่ต้องรองรับ” ที่แอปหลักมักตั้งไว้อยู่ระหว่าง Android 8 (API 26) ถึง Android 10 (API 29)—รองรับไปถึงเครื่องเมื่อไม่กี่ปีก่อน ฐานผู้ใช้ของคุณถึงจะกว้างพอ

\begin{center}
\begin{tabularx}{\linewidth}{|X|X|X|}
\hline
\textbf{เวอร์ชัน} & \textbf{ชื่อรหัส/ปี} & \textbf{สิ่งสำคัญสำหรับนักพัฒนา} \tabularnewline
\hline
Android 8 (API 26) & Oreo · 2017 & ช่องทางแจ้งเตือน ข้อจำกัดการทำงานเบื้องหลัง—เส้นต่ำสุดของเครื่องรุ่นเก่าจำนวนมาก \tabularnewline
Android 10 (API 29) & 2019 & เกิดระบบจัดเก็บแบบแบ่งส่วน (scoped storage) กฎการเข้าถึงไฟล์ปรับใหญ่ \tabularnewline
Android 12 (API 31) & 2021 & Material You ดึงสีแบบไดนามิก หน้าตาหน้าจออัปเกรดครั้งใหญ่ \tabularnewline
Android 13 (API 33) & 2022 & สิทธิ์แจ้งเตือนแบบขอตอนรันไทม์ มีป็อปอัปถามว่า "จะอนุญาตให้ส่งแจ้งเตือนไหม" \tabularnewline
Android 14-15 (API 34-35) & 2023-2024 & การย้อนกลับแบบคาดการณ์ หน้า 16KB ความเป็นส่วนตัวเข้มขึ้นอีกขั้น \tabularnewline
Android 16-17 (API 36-37) & 2025-2026 & จังหวะอัปเดตเวอร์ชันย่อยถี่ขึ้น บูรณาการความสามารถ AI (เช่น ส่วนขยาย Gemini) \tabularnewline
\hline
\end{tabularx}
\end{center}

\begin{center}\small\color{gray}ประวัติย่อเวอร์ชัน Android (ณ สิงหาคม 2026 สิ่งที่นักพัฒนาต้องโฟกัส)\end{center}

\booksection{เขียนด้วยภาษาอะไร: Java หรือ Kotlin?}

ประวัติศาสตร์การพัฒนา Android เคยมีสองภาษาหลัก Java คือ “ผู้ก่อตั้งรุ่นแรก”—Android เกิดขึ้นก็ใช้มันเลย; Kotlin คือ “พลเมืองชั้นหนึ่ง” ที่ Google สถาปนาปี 2017 และตั้งแต่ปี 2019 ทางการประกาศชัดว่า “Kotlin ก่อน” (Kotlin-first) ถึงปี 2026 โปรเจกต์ใหม่แทบทั้งหมดใช้ Kotlin โปรเจกต์เก่าส่วนใหญ่ก็กำลังย้ายไปเรื่อย ๆ

พูดกันตรง ๆ: ถ้าคุณเริ่มจากศูนย์จริง ๆ เรียน Kotlin ไปตรง ๆ อย่าย้อนไปยุ่งกับ Java เหตุผลมีสาม: หนึ่ง โค้ด Kotlin สั้นกว่า Java 30\%-40\% (ฟังก์ชันเดียวกัน Kotlin กระชับและชัดกว่า); สอง มันกัน “null pointer” ฝันร้ายของมือใหม่ได้ตั้งแต่ออกแบบ; สาม บทเรียนสมัยใหม่ทั้งหมด เอกสารทางการ และเครื่องมือเขียนโปรแกรม AI (ยังจำหนังสือเขียนโปรแกรม AI ของเราได้ไหม?) ใช้ Kotlin เป็นมาตรฐาน Java จะเจอแค่ตอนดูแลโปรเจกต์เก่าหรืออ่านโค้ดเก่าเท่านั้น

สไตล์ Java (เวอร์ชันเยอะคำ)

\begin{lstlisting}
public class MainActivity extends Activity \{
@Override
protected void onCreate(Bundle b) \{
super.onCreate(b);
setContentView(R.layout.activity);
\}
\}
\end{lstlisting}

สไตล์ Kotlin (เวอร์ชันกระชับ)

\begin{lstlisting}
class MainActivity : ComponentActivity() \{
override fun onCreate(b: Bundle?) \{
super.onCreate(b)
setContent \{ App() \}
\}
\}
\end{lstlisting}

ทำงานเดียวกัน Kotlin ตัดวงเล็บปีกกาและโค้ดบังคับไปเยอะ บทที่ 50 เราจะเรียนมันอย่างจริงจัง ตอนนี้คุณแค่ต้องจำว่า: “ภาษากลาง” ของการพัฒนา Android ชื่อว่า Kotlin

\booksection{แอปหนึ่งตัวถือกำเนิดขึ้นได้อย่างไร?}

หลายคนคิดว่า “พัฒนาแอป” ก็คือเปิดคอมพิวเตอร์แล้วพิมพ์โค้ดเละเทะ จริง ๆ แล้ว จากไอเดียจนถึงโทรศัพท์ผู้ใช้ แอปหนึ่งตัวต้องผ่านสายการผลิตประมาณนี้:

ไอเดีย

ออกแบบ (UI+ฟังก์ชัน)

เขียนโค้ด

คอมไพล์และแพ็ก (APK/AAB)

ทดสอบ (โปรแกรมจำลอง+เครื่องจริง)

เผยแพร่ (สโตร์/ไฟล์ติดตั้ง)

\begin{center}\small\color{gray}ภาพที่ 1: ขั้นตอนครบวงจรจากไอเดียถึงการขึ้นสโตร์ของแอป Android\end{center}

ในกระบวนการนี้มีสองคำที่คุณต้องรู้จักตั้งแต่ตอนนี้:

\begin{itemize}
  \item APK: รูปแบบไฟล์ติดตั้งของแอป Android (.apk) ดับเบิลคลิกก็ติดตั้งลงมือถือได้ เหมือน .exe ของ Windows ใช้ตอนทดสอบเอง หรือส่งให้เพื่อนเล่น
  \item AAB: Android App Bundle (.aab) รูปแบบบรรจุ "ขั้นสูง" ที่ใช้ขึ้น Google Play สโตร์จะตัดส่วนที่ไม่จำเป็นแบบไดนามิกตามรุ่นโทรศัพท์ของผู้ใช้ ทำให้ทุกคนดาวน์โหลดแพ็กเกจเล็กลง ติดตั้งเองไม่ได้ ต้องส่งให้สโตร์เท่านั้น
\end{itemize}

อีกหนึ่งคำที่คุณจะได้ยินบ่อย ๆ คือ SDK (Software Development Kit) มันคือชุด “อุปกรณ์พัฒนา” ครบเครื่องจาก Android ทางการ—รวมไลบรารี API ทุกเวอร์ชัน เครื่องมือดีบั๊ก อิมเมจโปรแกรมจำลอง ฯลฯ คลาสอย่าง Activity, Button ที่คุณเรียกใช้ตอนเขียนโค้ด ล้วนมาจาก SDK

\booksection{ภาพรวมเทคสแตกของ Android: “รายการอาวุธ” ของคุณ}

ก่อนเรียนอย่างเป็นทางการ มาทำความรู้จักการตั้งค่ามาตรฐานของโปรเจกต์ Android ยุคหลักปี 2026 ก่อน—ตอนนี้ดูรายการนี้แล้วไม่เข้าใจไม่เป็นไร พอเรียนจบทั้งเล่มกลับมาดูอีกที จะเข้าใจโปร่งเลย:

\begin{center}
\begin{tabularx}{\linewidth}{|X|X|X|}
\hline
\textbf{ชั้น} & \textbf{ทางเลือกหลัก (2026)} & \textbf{ใช้ทำอะไร} \tabularnewline
\hline
ภาษา & Kotlin 2.x & ภาษาที่ใช้เขียนตรรกะ \tabularnewline
อินเทอร์เฟซ & Jetpack Compose + Material 3 & วาดอินเทอร์เฟซด้วยโค้ด (สายหลักของเล่มนี้) \tabularnewline
สถาปัตยกรรม & MVVM + ViewModel & แยกอินเทอร์เฟซกับตรรกะออกจากกัน โค้ดไม่รก \tabularnewline
การจัดเก็บข้อมูล & Room (ฐานข้อมูล) + DataStore (การตั้งค่า) & เก็บข้อมูลลงโทรศัพท์ \tabularnewline
เครือข่าย & Retrofit + OkHttp + Kotlin coroutines & คุยกับเซิร์ฟเวอร์ \tabularnewline
เครื่องมือบิลด์ & Gradle + Kotlin DSL & เปลี่ยนโค้ดเป็นแพ็กเกจติดตั้ง \tabularnewline
สภาพแวดล้อมการพัฒนา & Android Studio & IDE "ครบเครื่อง" จากทางการ Google \tabularnewline
\hline
\end{tabularx}
\end{center}

\begin{center}\small\color{gray}สแตกเทคโนโลยี Android หลักปี 2026 ฉบับย่อ\end{center}

อย่าตกใจกับสายชื่อภาษาอังกฤษยาว ๆ พวกนี้ มันเหมือนหม้อหวดจานชามในครัว—ตอนนี้คุณแค่ต้องรู้ว่า “ทำกับข้าวต้องเข้าครัว” ส่วนเครื่องมือแต่ละชิ้นใช้ยังไง บทต่อ ๆ ไปจะแกะให้ดูทีละชิ้น และข่าวดีคือ: ตอน Android Studio สร้างโปรเจกต์ใหม่ มันจะตั้งค่าชุดนี้ให้เสร็จอัตโนมัติ คุณแค่เข้าครัวผัดเท่านั้น

\booksection{Android vs iOS: บทเรียนแรกของนักพัฒนาเรื่อง “เลือกฝั่ง”}

คุณอาจมีทั้งเครื่อง Android และ Apple อยู่ในมือ ในฐานะนักพัฒนา ความต่างของประสบการณ์สองฝั่งควรรู้ไว้ล่วงหน้า—ไม่ใช่การจับข้างว่า “ฝั่งไหนดีกว่า” แต่คือบัญชีความจริงว่า “ฝั่งไหนเหมาะกับการเริ่มต้นของคุณ” มากกว่า:

\begin{center}
\begin{tabularx}{\linewidth}{|X|X|X|}
\hline
\textbf{มิติ} & \textbf{พัฒนา Android} & \textbf{พัฒนา iOS} \tabularnewline
\hline
คอมพิวเตอร์ที่ใช้พัฒนา & Windows / Mac / Linux ก็ได้หมด & ต้อง Mac เท่านั้น (Xcode ใช้ได้เฉพาะ Mac) \tabularnewline
เครื่องมือพัฒนา & Android Studio (ฟรี) & Xcode (ฟรี) + ค่าบอกรายปี 99 ดอลลาร์ \tabularnewline
การตรวจสอบก่อนขึ้นสโตร์ & หลวม ใช้เวลาไม่กี่ชั่วโมงถึงไม่กี่วัน & เข้มงวด 1-7 วันคือเรื่องปกติ \tabularnewline
ติดตั้งลองเล่นเอง & APK ส่งได้ลงได้ตามใจ & การเผยแพร่มีข้อจำกัด การทดสอบต้องเซ็นแอป/ผ่าน TestFlight \tabularnewline
สภาพแวดล้อมการหารายได้ & พึ่งโฆษณา/ช่องทางในประเทศเป็นหลัก & ผู้ใช้ App Store จ่ายเงินซื้อของในแอปกันง่ายกว่า \tabularnewline
ต้นทุนการเรียนรู้ & มีคอนเซปต์เยอะ (เช่น คอมโพเนนต์หลักทั้งสี่) & อีโคไซสเต็มปิด แต่ระบบเป็นมาตรฐานเดียวกัน เอกสารมีคุณภาพ \tabularnewline
\hline
\end{tabularx}
\end{center}

\begin{center}\small\color{gray}เปรียบเทียบเส้นทางการพัฒนา Android กับ iOS (สิงหาคม 2026)\end{center}

พูดตรง ๆ กันเลยนะ: ถ้าคอมพิวเตอร์ของคุณไม่ใช่ Mac ก็ไม่ต้องลังเลอะไรเลย — Android คือทางเลือกเดียวที่มี แต่ถ้าคุณมี Mac ภาคที่ 8 กับ 9 ของหนังสือเล่มนี้เปิดเทียบกันอ่านได้ แนวคิดการเขียนโปรแกรมเหมือนกัน เรียนจบฝั่งหนึ่งแล้วไปเรียนอีกฝั่งจะไวขึ้นเยอะ

\begin{tcbpractice}{ลงมือทำด้วยกัน: ดู "Android" ในมือถือของคุณ}
เปิด [การตั้งค่า] บนมือถือ [เกี่ยวกับโทรศัพท์] (หรือ “เกี่ยวกับเครื่อง”) แล้วหาแถว “เวอร์ชัน Android” — แตะตรงนั้นติดกันหลายครั้ง ยังเปิดมินิเกมอีสเตอร์เอกซ์ของระบบได้อีกด้วย จดเลขเวอร์ชันของคุณไว้ แล้วเทียบกับตารางเวอร์ชันด้านบนดูว่าเป็นรุ่นที่เท่าไร มือถือของคุณนี่แหละคือ “เครื่องทดสอบ” ตัวแรกของคุณ

\end{tcbpractice}

\begin{tcbsummary}{สรุปบทนี้}
\begin{itemize}
  \item Android เกิดปี 2008 และปี 2026 เป็นระบบปฏิบัติการมือถือที่มีส่วนแบ่งตลาดโลกสูงที่สุด มีอุปกรณ์ใช้งานจริงหลายพันล้านเครื่อง
  \item เวอร์ชันทางการล่าสุดคือ Android 17; แอปยอดนิยมรองรับขั้นต่ำ Android 8-10 เพื่อให้ฐานผู้ใช้ใหญ่พอ
  \item โปรเจกต์ใหม่ใช้ Kotlin ไปเลยทั้งหมด: โค้ดน้อยลง 30\%+ กัน NullPointerException และเป็นภาษาที่ทางการรับรอง
  \item APK ติดตั้งเองได้ AAB ไว้ขึ้นสโตร์; SDK คือชุดเครื่องมือพัฒนาอย่างเป็นทางการ
  \item สแตกเทคโนโลยีมาตรฐานปี 2026: Kotlin + Compose + MVVM + Room + Retrofit + Gradle + Android Studio
\end{itemize}

\end{tcbsummary}
